% ===== 原文 pp.51–54（PDF p.61–64） =====

% [原文 p.51]
如果把 $A_4$ 换成它的和乐 $h$、把 $\rho$ 换成 $\hat\rho$（如我们对 Vafa-Witten 理论所做的那样）并约化到三维，那么所得的 3d 约化作用量与 Vafa-Witten 理论约化所得的 (5.3) 相符，只差项
\[
    \frac{1}{e^2}\delta_Q\Tr\{4\tilde\lambda^r[b_r,\phi]\}
\]
的符号，而这一项的符号改变会使 (5.4) 中 $[\star b,\phi]$ 项的符号改变。然而，BPS 方程的解反正会使该项为零，所以这个符号无关紧要，我们由此断定 Vafa-Witten 与 $t=0$ 两种理论的约化是相同的。

% [Tutor Note: p.51 的“符号无关紧要”最好用另一种方式核对：两个约化作用量之差是某个 $\delta_Q$-正合项，按 §5.2 的形变论证它不改变拓扑关联函数；若直接比较 (5.3) 与 $t=0$ 的约化式，可看到差别只落在 $\cL_3$ 的 $4\tilde\lambda^r[b_r,\phi]$ 一项，其 $\delta_Q$ 作用后的动能项符号随之改变，故应以 $\delta_Q$-正合性（而非“该项为零”）为判定依据。]

\section{GL-扭转理论在 \texorpdfstring{$t=1$}{t=1} 时的 \texorpdfstring{$S^1$}{S1} 紧化}

$t=1$ 理论的约化同样给出 3d A-模型的带规范场版本。然而，与 Vafa-Witten 及 $t=0$ 的约化相比，有几个重要的差别。场 $A_4$ 现在将位于 3d 理论的物质扇区（而不是规范扇区）；确实，约化理论的靶空间将是群 $G$，由和乐 $h$ 参数化。相应地，物质扇区的场取值于向量空间 $T_hG$ 而不是 $\fg\simeq T_1G$。此外，作用量中拓扑项的系数只是 $\theta$ 而不是复化耦合 $\tau$，这意味着我们可以通过令 $\theta=0$ 来避开解释瞬子项的麻烦。最后，规范扇区将是一个通常的 3d A-型规范理论，如 B.2.1 节所述（而不是涉及群取值标量的奇特版本）。

% [Tutor Note: p.51 的三点差别可逐条对应到后面的公式：$A_4$ 属物质扇区体现为 p.52 物质表里的 $\phi=A_4^{[1]}$；靶空间为 $G$ 体现为 p.52 的 $\Phi=h:W\to G$；拓扑项系数为 $\theta$ 体现为 p.51 作用量中的 $+\ii\theta/(8\pi^2)$（对比 (5.6) 的 $\bar\tau$、§5.2 的 $\tau$）。]

[1] 中 GL-扭转理论在 $t=1$ 时的作用量为
\[
    S = \int_M\dd^4x\sqrt{g}\left(\cL_1+\cL_2\right)
    + \frac{\ii\theta}{8\pi^2}\int_M\Tr F\wedge F
\]
其中
\begin{align}
    \cL_1 &= \frac{1}{e^2}\delta_Q\Tr\left\{\chi_{\mu\nu}
        \left(\half H - F + \half[\phi,\phi] - \ast\dd_A\phi\right)^{\mu\nu}\right\} \notag\\
    \cL_2 &= \frac{1}{e^2}\delta_Q\Tr\left\{-\tilde\eta\left(P - 2\dd_A^\star\phi\right)
        + \bar\sigma\left(2\dd_A^\star\psi + 2[\tilde\psi_\mu,\phi^\mu] + [\sigma,\eta]\right)\right\} . \tag{5.7}
\end{align}

重新标记场后，规范扇区的场为
\begin{gather*}
    A_r = A_r^{[1]}, \qquad \lambda_r = \ii\left(\psi_r^{[1]} + \tilde\psi_r^{[1]}\right) \\
    \varsigma = \phi_4^{[1]}, \qquad \rho = \ii\left(\psi_4^{[1]} - \tilde\psi_4^{[1]}\right) \\
    \sigma = -2\ii\sigma^{[1]}, \\
    \bar\sigma = -2\ii\bar\sigma^{[1]}, \qquad \tilde\rho = 2\left(\eta^{[1]} + \tilde\eta^{[1]}\right), \\
    \tilde\lambda_r = (\ast\chi)_{r4}^{[1]}, \qquad H_r = (\ast H)_{r4}^{[1]},
\end{gather*}

% [原文 p.52]
物质扇区的场为
\begin{gather*}
    \phi = A_4^{[1]}, \qquad \eta = \ii\left(\psi_4^{[1]} + \tilde\psi_4^{[1]}\right), \\
    b_r = \phi_r^{[1]}, \qquad \psi_r = \ii\left(\psi_r^{[1]} - \tilde\psi_r^{[1]}\right), \\
    \beta = \half\left(\eta^{[1]} - \tilde\eta^{[1]}\right), \qquad P = P^{[1]}, \\
    \chi_r = \chi_{r4}^{[1]}, \qquad \tilde P_r = H_{r4}^{[1]} .
\end{gather*}
重新标记后的场的 BRST 变分再一次由 (5.1) 给出。把 $t=1$ 的作用量用重新标记后的场写出，我们得到
\[
    S = \int_M\dd^4x\sqrt{g}\left(\cL_1+\cL_2+\cL_3+\cL_4+\cL_5\right)
\]
其中
\begin{align}
    \cL_1 &= \frac{1}{e^2}\delta_Q\Tr\left\{\tilde\lambda^r\left(H_r - 2D_r\varsigma - 2\star F_r\right) + \half\bar\sigma\, \dd_A^\star\lambda\right\} \notag\\
    \cL_2 &= \frac{1}{e^2}\delta_Q\Tr\left\{\chi^r\left(\tilde P_r - 2D_rA_4 - 2(\star \dd_Ab)_r\right) + \beta\left(P - 2\dd_A^\star b\right)\right\} \notag\\
    \cL_3 &= \frac{1}{e^2}\delta_Q\Tr\left\{\tilde\lambda^r\left(2D_4b_r + \star[b,b]_r\right) + \chi^r\left(2[b_r,\varsigma] + 2\partial_4A_r\right)\right\} \notag\\
    \cL_4 &= \frac{1}{e^2}\delta_Q\Tr\left\{2\beta\left(-D_4\varsigma - \frac{1}{8}[\bar\sigma,\sigma]\right)
        - \half\tilde\rho\left(\half P - \dd_A^\star b - D_4\varsigma + \frac{1}{8}[\bar\sigma,\sigma]\right)\right\} \notag\\
    \cL_5 &= \frac{1}{e^2}\delta_Q\Tr\left\{\half\bar\sigma\left(\dd_A^\star\psi + D_4\eta + D_4\rho + [\lambda_r,b^r] - [\psi_r,b^r] + [\eta,\varsigma] - [\rho,\varsigma]\right)\right\} . \tag{5.8}
\end{align}

% [Tutor Note: p.52 的场表可按 (5.1) 逐条“字典式”核对：例如由 (5.1) 的 $\delta_QA=\lambda$ 与 p.51 的 $\lambda_r=\ii(\psi_r^{[1]}+\tilde\psi_r^{[1]})$ 可反推 $t=1$ 的 BRST 变分必须是 $\delta_QA_r=\ii(\psi_r+\tilde\psi_r)$；再由 $\delta_Q\varsigma=\rho$ 得 $\delta_Q\phi_4=\ii(\psi_4-\tilde\psi_4)$，与 (5.8) 中 $\cL_2$ 的 $D_rA_4$ 项一致。若某一行归一化写错（如漏掉 $\sqrt2$ 或 $\ii$），这条字典会立刻矛盾。]

我们再把理论描述中的 $A_4$ 换成它的和乐 $h$。靶空间是群流形 $X=G$，$\sigma$-模型映射 $\Phi$ 由和乐给出：
\[
    \Phi = h : W \to G .
\]
3d A-模型场 $\eta,b_r,\beta,P,\chi_r,\tilde P_r$ 现在取值于 $T_hG$ 而不是李代数本身；把它们换成 $\fg$ 取值的平均值是方便的，正如我们对上述 Vafa-Witten 约化中的 $\rho$ 所做的那样，即
\[
    \hat\eta \equiv -\frac{1}{2\pi}\int_0^{2\pi}\dd x^4\, H^{-1}\eta H
\]
对其余物质场类似。平均后场的 BRST 变分多出一项（来自群流形的曲率），如下：
\[
    \delta_Q\hat b_r = \hat\psi_r + \langle\hat\eta,\hat b_r\rangle, \qquad
    \delta_Q\hat\chi_r = \hat P_r + \langle\hat\eta,\hat\chi_r\rangle
\]

% [原文 p.53]
其中我们定义了
\begin{align*}
    \langle\hat\eta,\hat b_r\rangle &\equiv \int_0^{2\pi}\dd x^4
        \left[\hat\eta,\, H^{-1}b_rH\right]\Big|_{x^4}, \\
    \langle\hat\eta,\hat\chi_r\rangle &\equiv \int_0^{2\pi}\dd x^4
        \left\{\hat\eta,\, H^{-1}\chi_rH\right\}\Big|_{x^4} .
\end{align*}

% [Tutor Note: p.53 两个括号一为对易子 $[\ ,\ ]$、一为反对易子 $\{\ ,\ \}$，原因是 $b_r$ 为玻色场而 $\chi_r$ 为费米场；抄录时若把 $\{,\}$ 写成 $[,]$ 会破坏 $\delta_Q^2=0$（因为 $\hat\eta$ 是费米场）。这是检验此处符号的最快方法。]
% [Tutor Note: p.53 的附加项 $\langle\hat\eta,\hat b_r\rangle$ 正是群流形上的“联络项”：$T_hG$ 取值的场在平均后要用 $\fg$ 取值的场表示，须计入 $H^{-1}\dd H$ 型的和乐修正；它与 §5.2 中 $h^{-1}D_rh$ 的出现同源，可对照 (5.3) 的 $\hat\cL_1$ 理解。]

约化步骤将与我们上述约化 Vafa-Witten 作用量所描述的类似：我们用一个 BRST-正合的量光滑地形变作用量，以便用 $h$ 以及 $T_hG$ 取值物质场的 $\fg$ 取值平均值来取代 $A_4$ 的出现。约化到 3d 后，得到
\[
    \hat S = \int_W\dd^3x\sqrt{g_W}\left(\hat\cL_1+\hat\cL_2+\hat\cL_3+\hat\cL_4+\hat\cL_5\right)
\]
其中
\begin{align*}
    \hat\cL_1 &= \frac{1}{e^2}\delta_Q\Tr\left\{\tilde\lambda^r\left(H_r - 2D_r\varsigma - 2\star F_r\right) + \half\bar\sigma\, \dd_A^\star\lambda\right\} \\
    \hat\cL_2 &= \frac{1}{e^2}\delta_Q\Tr\left\{\hat\chi^r\left(\hat P_r + 2(2\pi)^{-1}h^{-1}D_rh - 2(\star \dd_A\hat b)_r\right)
        + \hat\beta\left(\hat P - 2\dd_A^\star\hat b\right)\right\} \\
    \hat\cL_3 &= \frac{1}{e^2}\delta_Q\Tr\left\{\tilde\lambda^r\left(-2(2\pi)^{-1}h^{-1}[h,\hat b_r] + \star[\hat b,\hat b]_r\right)
        + 2\hat\chi^r[\hat b_r,\varsigma]\right\} \\
    \hat\cL_4 &= \frac{1}{e^2}\delta_Q\Tr\left\{2\hat\beta\left((2\pi)^{-1}h^{-1}[h,\varsigma] - \frac{1}{8}[\bar\sigma,\sigma]\right)
        - \half\tilde\rho\left(\half\hat P - \dd_A^\star\hat b + (2\pi)^{-1}h^{-1}[h,\varsigma] + \frac{1}{8}[\bar\sigma,\sigma]\right)\right\} \\
    \hat\cL_5 &= \frac{1}{e^2}\delta_Q\Tr\left\{\half\bar\sigma\left(\dd_A^\star\hat\psi - (2\pi)^{-1}h^{-1}[h,\hat\eta]
        - (2\pi)^{-1}h^{-1}[h,\rho] + [\lambda_r,\hat b^r] - [\hat\psi_r,\hat b^r] + [\hat\eta,\varsigma] - [\rho,\varsigma]\right)\right\} .
\end{align*}
把辅助场积掉后，人们发现理论局域化在 BPS 方程
\begin{align*}
    F + \star \dd_A\varsigma - \half[\hat b,\hat b] + (2\pi)^{-1}h^{-1}[h,\star\hat b] &= 0 \\
    \star \dd_A\hat b - [\hat b,\varsigma] - (2\pi)^{-1}h^{-1}D_rh &= 0 \\
    \dd_A^\star\hat b - (2\pi)^{-1}h^{-1}[h,\varsigma] &= 0 .
\end{align*}
上。

我们得到了靶空间为 $G$ 的带规范场的 3d A-模型。这个 3d TFT 在文献中研究得不多，但由于它被猜测与 $t=\ii$ 理论的 3d 约化对偶，研究其性质预期会带来许多有趣的数学预言。在下一章中，我们将通过在 $G=U(1)$ 情形下分析它的边界条件的 2-范畴，朝这个方向迈出不大的第一步。

% [原文 p.54]
\section{GL-扭转理论在 \texorpdfstring{$t=\ii$}{t=i} 时的 \texorpdfstring{$S^1$}{S1} 紧化}

$t=\ii$ 时 GL-扭转理论的约化可以用类似的方式分析；然而，所得的 3d 约化理论将具有与迄今考虑的情形非常不同的特征。人们得到的是一个被充分研究过的 3d TFT——Rozansky-Witten（RW）模型 [14], [36]——的带规范场版本。（带规范场是通过耦合到 B.4 节所述的 3d B-型规范理论来实现的。）该模型的靶空间是复化李群的余切丛，$X=T^*G_{\C}$，并带有 $G$ 在底空间 $G_{\C}$ 上的共轭作用以及纤维上诱导的作用。这来自如下机制。

$t=\ii$ 理论中 BRST 不变的复场 $\sigma$ 与 $A_4+\ii\phi_4$ 将下传到 3d 理论的物质扇区。如我们所讨论的，场 $A_4$ 本身定义不良，必须换成一个适当的群取值标量。我们定义
\[
    h_{\C} = P\exp\left(-\int_0^{2\pi}\dd x^4\left(A_4+\ii\phi_4\right)\right)
\]
取值于纤维为 $G_{\C}$ 的相伴丛：
\[
    h_{\C} \in \Omega^0(W, P_W\times_G G_{\C}) .
\]
其共轭为
\[
    \bar h_{\C} = P\exp\left(-\int_0^{2\pi}\dd x^4\left(A_4-\ii\phi_4\right)\right) .
\]
场 $h_{\C}$ 是底空间的好坐标，而 $\sigma$ 是 $T^*G_{\C}$ 的纤维的好坐标。流形 $T^*G_{\C}$ 被证明带有超 Kähler 度规 [63]，尽管显式描述它并不简单。

% [Tutor Note: p.54 的 $h_{\C}\in\Omega^0(W,P_W\times_GG_{\C})$ 与 $h$ 的情形同理：复化联络 $A+\ii\phi$ 的绕圈和乐取值于 $G_{\C}$，规范变换按共轭作用，故它是纤维为 $G_{\C}$ 的相伴丛的截面。核对维数：$T^*G_{\C}$ 的实维数为 $2\dim G$，底坐标 $h_{\C}$ 贡献 $\dim G$、纤维坐标 $\sigma\in\fg_{\C}$ 贡献 $\dim G$，正好吻合。]
% [Tutor Note: p.54 的 $\Z_2\to U(1)$ 提升是关键约定：一般超 Kähler 靶空间的 RW 模型只有 $\Z_2$ 鬼数对称性（$\sigma$ 的鬼数只能是 $2\bmod 2$），而余切丛的纤维坐标可整体赋予鬼数二，于是 $U(1)$ 鬼数可以定义——这也解释了记号 $T^*[2]G_{\C}$ 中 $[2]$ 的含义（鬼数移位）。]

简单地把靶空间宣布为 $T^*G_{\C}$ 会遮蔽作用量的一个重要对称性。对一般的超 Kähler 靶空间 $X$，RW 模型具有 $\Z_2$ 鬼数对称性；但如 [36] 所解释的，当 $X$ 是余切丛时，可以令纤维坐标具有鬼数二，从而把它提升为 $U(1)$ 鬼数对称性。这与 $\sigma$ 在 4d 理论中已经具有鬼数二这一事实很好地吻合。因此，为强调纤维坐标被形式地赋予鬼数二，我们将改用 $T^*[2]G_{\C}$ 记靶流形。

$t=\ii$ 时 GL-扭转理论的作用量写在式 (2.4) 中。场的重新标记如下（采用 $[i]$ 上标表示 $t=\ii$ 理论中场的约定，且
